\section{Related Work}
\label{sec:related_work}

The literature on model combination, adaptation, and weight-space dynamics forms the foundation for our deconstruction of test-time model merging. We categorize the closely related work into three primary domains: static model merging, test-time adaptation, and representational regularizations.

\subsection{Static Weight-Space Model Merging}
Model merging is the practice of combining multiple fine-tuned models sharing a common pre-trained ancestor without accessing their training datasets. Early work on Model Soups \cite{wortsman2022model} demonstrated that averaging the weights of models fine-tuned with different hyperparameters on the same task improves generalization and calibration. This concept was extended to multi-task settings via Task Arithmetic \cite{ilharco2022editing}, which adds task-specific "task vectors" (the difference between fine-tuned and pre-trained weights) to the pre-trained base model. 

However, naive task vector addition can lead to severe destructive interference when parameters conflict in sign or magnitude. To mitigate this, TIES-Merging \cite{yadav2023resolving} introduced a spatial consensus framework that filters out small parameter changes, resolves sign conflicts, and averages only the agreeing parameters. Similarly, DARE \cite{yu2023language} uses randomized pruning and scaling to merge massive language models. While these static approaches are parameter-efficient and training-free, they employ a uniform coefficient across all layers and tasks, which fails to account for the highly localized and variable nature of multi-task representations across deep network hierarchies.

\subsection{Adaptive and Test-Time Model Merging}
To overcome the limitations of uniform, static merging, adaptive merging frameworks optimize merging coefficients across different layers or modules. RegMean \cite{jin2023regmean} optimizes weights by minimizing the difference in activations using cached training data, while Fisher-weighted merging \cite{matena2022merging} uses Fisher information matrices to weigh parameters. However, these methods require either training-set activations or expensive Hessian estimations.

Test-Time Adaptation (TTA) addresses this by optimizing parameters directly on unlabeled test streams. Grounded in unsupervised objectives like entropy minimization (e.g., Tent \cite{wang2021tent}), TTA adapts model parameters on the fly to handle distribution shifts. This test-time adaptation paradigm was applied to model merging by AdaMerging \cite{yang2024adamerging}, which optimizes layer-wise merging coefficients to minimize joint prediction entropy on unlabeled calibration streams. SyMerge \cite{jung2025symerge} further extended this by introducing single-layer adaptation mechanisms to induce positive task synergy. Yet, as we demonstrate, these unconstrained test-time optimizers operate on highly restricted calibration batches (often as small as 16 samples) and suffer from severe transductive overfitting and task-sacrifice biases, which have been largely ignored in prior literature.

\subsection{Linear Mode Connectivity and Representation Alignment}
The success of model merging is deeply linked to Linear Mode Connectivity (LMC), which states that models fine-tuned from a shared pre-trained initialization lie within the same low-loss basin and can be linearly interpolated without encountering high-energy barriers \cite{frankle2020linear, mirzadeh2020linear}. When models are trained from different initializations, they occupy disparate basins, requiring permutation alignment techniques like Git Re-Basin \cite{ainsworth2022git}, ZipIt! \cite{zhang2023zipit}, or representation repair mechanisms like REPAIR \cite{jordan2022repair} to re-normalize activation paths and align representation manifolds.

In our work, we do not focus on aligning disparate basins; instead, we study the optimization surface of models already within a shared basin. We build on representational regularizations in neural networks by demonstrating that standard, fine-grained layer-wise parameter updates in test-time adaptation are heavily overfit. By introducing spatial constraints (e.g., penalizing spatial deviations from the task average), we bridge the gap between uniform static merging (which has zero spatial degrees of freedom but high robustness) and fully adaptive test-time merging (which has high spatial degrees of freedom but suffers from transductive overfitting). This provides a calibrated, smoothly adjustable framework for multi-task weight interpolation.